\section{はじめに}
%----------------------------------------------------------------------

近年，スマートフォンや医療機器，IoT センサなど多様な端末から生成されるデータを活用した
機械学習の需要が急速に高まっている．
しかしこれらのデータは高度に個人的であり，
プライバシーの観点から端末外へのデータ転送が困難または不可能なケースが多い．

こうした背景から，McMahan ら\cite{McMahanFedAvg}が提案した
連合学習（Federated Learning; FL）が注目を集めている．
FL はデータを端末内に保持したまま，各端末がローカルで学習したモデルパラメータのみをサーバと
共有することで，プライバシーを保護しつつ分散的な協調学習を実現する枠組みである．
医療診断，スマートフォンの次単語予測，スパム検出など幅広い応用が報告されている．

しかし実用的な FL 環境においては，2 種類の異種性が同時に存在するという問題がある．

\textbf{計算異種性（Computational Heterogeneity）}:
スマートフォン，タブレット，組み込みデバイスなど，
処理能力・メモリ容量が大きく異なる端末が混在する．
計算資源が限られた端末はフルサイズのモデルを実行できないため，
すべての端末が同一のモデルを用いる従来の FL 手法（FedAvg\cite{McMahanFedAvg}）では，
非力な端末が学習に参加できないという問題が生じる．

\textbf{アーキテクチャ異種性（Architectural Heterogeneity）}:
端末や組織ごとに異なるモデルアーキテクチャを採用するケースが増えている．
たとえば画像認識において，CNN ベースのモデルと
Vision Transformer（ViT）\cite{VisionTransformer}ベースのモデルが混在することは珍しくない．
異なるアーキテクチャ間では重みの構造が異なるため，
FedAvg のようなパラメータ平均による集約が直接適用できない．

既存の代表的手法はいずれか一方の異種性にしか対応できない（表~\ref{tab:related}）．
HeteroFL\cite{HeteroFL}は計算異種性を解決するが，
異なるアーキテクチャ間での知識共有を提供しない．
FedGen\cite{FedGen}はアーキテクチャ非依存の知識蒸留を実現するが，
全端末が同一サイズのモデルを持つことを前提としている．

本研究では，この空白を埋めるフレームワーク
AFAD（Adaptive Federated Architecture Distribution）を提案する．
AFAD は，32次元の共有ボトルネック層を介して HeteroFL と FedGen を統合することで，
計算異種性とアーキテクチャ異種性の両方を同時に扱う．
本稿の貢献を以下に示す．

\begin{itemize}
  \item 著者らの知る限り，計算異種性（モデル幅）とアーキテクチャ異種性（CNN/ViT）を
        同時に扱う最初の FL フレームワーク AFAD を提案する．
  \item 32次元共有ボトルネック層という単純な設計により，
        HeteroFL の重み集約機構と FedGen の知識蒸留機構を干渉を最小化して統合する設計を示す．
  \item OrganAMNIST を用いた IID 環境での実験により，
        AFAD がサーバ精度 82.93\% を達成し，HeteroFL（82.20\%）を上回ることを示す．
        二重異種性への同時対応が HeteroFL 比での精度低下を伴わないことを実証する．
\end{itemize}

%----------------------------------------------------------------------
